% الجزء: sets-functions-relations
% الفصل: relations
% القسم: special-properties

\documentclass[../../../include/open-logic-section]{subfiles}

\begin{document}

\olfileid[ar]{sfr}{rel}{prp}
\olsection{خصائص خاصة للعلاقات}

\begin{intro}
من العلاقات أصناف كثرت حتى اختصت بأسماء. فـ$\le$ و~$\subseteq$
مثلًا تتشابهان في ربط عناصر مجاليهما: $\Nat$ في حالة~$\le$،
و$\Pow{A}$ في حالة~$\subseteq$. ولتمييز موضع الاتفاق والافتراق
نصنف العلاقات بخصائص يمكن أن تتصف بها. وبعض تلك الخصائص، منفردة
أو مجتمعة، بالغ الأهمية؛ ومن اجتماعها تنشأ الترتيبات وعلاقات التكافؤ.
\end{intro}

\begin{defn}[الانعكاسية]
العلاقة $R \subseteq A^2$ \emph{انعكاسية} إذا وفقط إذا كان كل $x \in
A$ يحقق $Rxx$.
\end{defn}

\begin{defn}[التعدية]
العلاقة $R \subseteq A^2$ \emph{متعدية} إذا وفقط إذا لزم من $Rxy$
و$Ryz$ معًا $Rxz$.
\end{defn}

\begin{defn}[التناظر]
العلاقة~$R \subseteq A^2$ \emph{متناظرة} إذا وفقط إذا لزم من
$Rxy$ صدق~$Ryx$.
\end{defn}

\begin{defn}[مضادّة التناظر]
العلاقة~$R \subseteq A^2$ \emph{مضادّة\- للتناظر\-} إذا وفقط إذا لم
يجتمع $Rxy$ و$Ryx$ إلا مع $x=y$؛ أي متى كان $x\neq y$ صدق أحد
$\lnot Rxy$ و$\lnot Ryx$ على الأقل.
\end{defn}

\begin{explain}
أما في المتناظرة فـ$Rxy$ و$Ryx$ إما صادقان معًا وإما كاذبان معًا.
وأما في المضادّة للتناظر فلا يجتمع $Rxy$ و$Ryx$ إلا إذا كان
$x = y$. وليس في ذلك \emph{اشتراط} صدق $Rxy$ و$Ryx$ عند $x = y$،
وإنما فيه جوازه. لذلك قد تكون المضادّة للتناظر انعكاسية، ولا يلزم
أن تكون كل مضادّة للتناظر كذلك. ثم لا تخلط مضادّة التناظر بمجرد
انتفاء التناظر؛ فهما شرطان مختلفان. وقد يجتمع التناظر ومضادّته
في علاقة واحدة، كما في علاقة الهوية.
\end{explain}

\begin{defn}[الاتصالية]
العلاقة $R \subseteq A^2$ \emph{متصلة} متى كان لكل $x,y\in
A$، مع $x \neq y$، أحد $Rxy$ و~$Ryx$ صادقًا على الأقل.
\end{defn}

\begin{prob}
مثّل بعلاقات تحقق هذه الأوصاف: (أ) الانعكاسية والتناظر دون التعدية؛
(ب) الانعكاسية ومضادّة التناظر؛ (ج) مضادّة التناظر والتعدية دون
الانعكاسية؛ (د) الانعكاسية والتناظر والتعدية جميعًا. وليست الأمثلة
المطلوبة علاقات على الأعداد ولا على المجموعات.
\end{prob}

\begin{defn}[اللاانعكاسية]
العلاقة $R \subseteq A^2$ \emph{لا انعكاسية} متى امتنع لكل $x \in
A$ أن يصدق $Rxx$.
\end{defn}

\begin{defn}[اللاتناظرية]
العلاقة $R \subseteq A^2$ \emph{لا تناظرية} متى امتنع وجود زوج
$x,y\in A$ يجتمع عنده $Rxy$ و~$Ryx$.
\end{defn}

فإذا كان $A \neq \emptyset$ لم تجتمع اللاانعكاسية والانعكاسية في
علاقة على~$A$، وكانت كل لا تناظرية على~$A$ مضادّة للتناظر.
غير أنه توجد علاقات $R \subseteq A^2$ ليست انعكاسية ولا لا انعكاسية؛
وتوجد كذلك مضادّات للتناظر ليست لا تناظرية.

\end{document}
